\section{Integrals}

An integral describes accumulation. It can be understood as a limiting sum or, geometrically, as the signed area under a graph. In this introductory chapter we mainly need the connection between integration and antiderivatives.

\begin{definitionbox}
Let \(f\colon I\to\mathbb{R}\) be a function defined on an interval \(I\). A function \(F\colon I\to\mathbb{R}\) is called an antiderivative of \(f\) on \(I\) if \(F\) is differentiable on \(I\) and
\[
F'(x)=f(x)
\]
for every \(x\in I\).
\end{definitionbox}

\begin{definitionbox}
The indefinite integral of \(f\) is the family of all antiderivatives of \(f\). If \(F\) is one antiderivative of \(f\), we write
\[
\int f(x)\,dx=F(x)+C,
\]
where \(C\) is an arbitrary constant.
\end{definitionbox}

For example,
\[
\int 2x\,dx=x^2+C,
\]
because
\[
\frac{d}{dx}(x^2+C)=2x.
\]

\begin{remarkbox}
The constant \(C\) appears because the derivative of a constant is zero. Therefore many functions can have the same derivative.
\end{remarkbox}

The definite integral of \(f(x)\) from \(a\) to \(b\) is written as
\[
\int_a^b f(x)\,dx.
\]
It measures accumulation over the interval \([a,b]\). Its precise construction can be given using limiting sums; in these notes we will mainly compute it using the theorem below.

\begin{theorembox}
If \(f\) is continuous on \([a,b]\) and \(F\) is an antiderivative of \(f\) on \([a,b]\), then
\[
\int_a^b f(x)\,dx=F(b)-F(a).
\]
\end{theorembox}

\begin{examplebox}
Since \(F(x)=x^2\) is an antiderivative of \(f(x)=2x\), we have
\[
\int_0^3 2x\,dx
=
\left[x^2\right]_0^3
=
9-0
=
9.
\]
\end{examplebox}

The symbol \(dx\) indicates the variable of integration. In the expression
\[
\int_a^b f(x)\,dx,
\]
we integrate with respect to \(x\).

Some basic integration rules are
\[
\int x^n\,dx=\frac{x^{n+1}}{n+1}+C \quad (n\neq -1),
\]
\[
\int (f(x)+g(x))\,dx=\int f(x)\,dx+\int g(x)\,dx,
\qquad
\int c f(x)\,dx=c\int f(x)\,dx.
\]

\begin{examplebox}
\[
\int (3x^2+5x-7)\,dx
=
x^3+\frac{5}{2}x^2-7x+C.
\]
\end{examplebox}

In this chapter, the important mathematical idea is that integration describes accumulation and is closely connected with differentiation.
